\begin{figure}[!htbp]
\centering
\resizebox{0.98\linewidth}{!}{%
\begin{tikzpicture}[x=1cm,y=1cm,>=Latex,font=\small,
  timeaxis/.style={->,black!65,line width=.55pt},
  hittick/.style={blue!55!black,line width=1.1pt},
  horizontick/.style={orange!75!black,line width=1.1pt}]
  \path[use as bounding box] (0,-.32) rectangle (13.7,5.97);
  \node[anchor=west,align=left] at (0,5.32)
    {Physical time\\$t$};
  \node[anchor=west,align=left] at (0,3.33)
    {Hit dilation\\$t/r_N$};
  \node[anchor=west,align=left] at (0,1.33)
    {Horizon dilation\\$t/T_N$};

  \draw[black!20,dashed] (2.45,.99)--(2.45,5.37);
  \draw[blue!35!black,dashed] (6.70,.99)--(6.70,5.37);
  \draw[orange!55!black,dashed] (12.05,.99)--(12.05,5.37);
  \foreach \yy in {5.30,3.30,1.30}{
    \draw[timeaxis] (2.20,\yy)--(13.15,\yy);
    \draw[black!60] (2.45,\yy-.10)--(2.45,\yy+.10);
    \draw[hittick] (6.70,\yy-.13)--(6.70,\yy+.13);
    \draw[horizontick] (12.05,\yy-.13)--(12.05,\yy+.13);
    \node[below,fill=white,inner sep=2pt] at (2.45,\yy-.14) {$0$};
  }
  \node[above,text=blue!55!black] at (6.70,5.51) {cap hit};
  \node[above,text=orange!75!black] at (12.05,5.51) {horizon};
  \node[below,fill=white,inner sep=2pt] at (6.70,5.12)
    {$t_{*,N}=r_N-s_N$};
  \node[below,fill=white,inner sep=2pt] at (12.05,5.12) {$T_N$};
  \node[below,fill=white,inner sep=2pt] at (6.70,3.12) {$1-l_N$};
  \node[below,fill=white,inner sep=2pt] at (12.05,3.12) {$U_N=T_N/r_N$};
  \node[below,fill=white,inner sep=2pt] at (6.70,1.12)
    {$\dfrac{1-l_N}{U_N}$};
  \node[below,fill=white,inner sep=2pt] at (12.05,1.12) {$1$};

  \node[align=center,fill=white,inner sep=3pt] at (7.25,4.25)
    {$r_N=t_{*,N}+s_N,\qquad l_N=s_N/r_N$};
  \node[align=center,fill=white,inner sep=3pt] at (7.25,2.24)
    {$t/T_N=(t/r_N)/U_N,\qquad x/T_N=(x/r_N)/U_N$};
  \node[align=center] at (7.05,-.12)
    {$U_N\longrightarrow U<\infty
      \quad\Longrightarrow\quad r_N/T_N\longrightarrow 1/U>0$};
\end{tikzpicture}%
}
\caption{The same hit and horizon in the two dilations.
The rows align physical events and are schematic; the hit may occur at either
endpoint, and no ordering of $r_N$ and $T_N$ is assumed.
Sections~\ref{sec:dilation}--\ref{sec:branching} will show that the hit
dilation has a nonzero weak limit on $0<t<U$ and that $U<\infty$.
The latter conclusion makes the change to the horizon dilation nondegenerate:
compact subsets of $0<t<U$ map inside $0<t<1$.
Thus vanishing in the horizon dilation would force vanishing in the hit
dilation, which is the contradiction used to separate the solutions.}
\label{fig:scales}
\end{figure}
